\documentclass[11pt,a4paper]{article}
\usepackage[margin=2.5cm]{geometry}
\usepackage{parskip}
\usepackage[T1]{fontenc}
\usepackage{lmodern}
\usepackage{titlesec}
\titleformat{\section}{\large\bfseries}{\thesection}{0.6em}{}

\title{Implementation Explanations\\[0.4em]\large The Enrichment Analysis (Pipeline Step 10)}
\author{MH-Neuron / Modality Taxonomy project}
\date{June 10, 2026}

\begin{document}
\maketitle

The enrichment analysis is the experiment that tests whether the neuron labels the
pipeline produces actually mean something causally, not just statistically. Here is
the whole thing in plain words, start to finish.

\section*{The starting point}

By this stage of the pipeline, every neuron in the language model's feed-forward
layers has already been given a label --- visual, text, multimodal, or unknown ---
based on \emph{when it tends to fire}: neurons that activate mostly on image-derived
tokens are called visual, those that activate mostly on language tokens are called
text, and so on. But firing patterns are only correlations. A neuron that fires on
visual tokens might still contribute nothing to the model's visual understanding.
The enrichment analysis asks the causal question: if we actually \emph{remove} these
neurons, does the model's behavior break in the way the labels predict?

\section*{The behavioral yardstick}

The behavior chosen is hallucination --- specifically, object hallucination measured
by POPE, a benchmark of simple yes-or-no questions like ``Is there a dog in this
image?'' asked about real photographs. A model hallucinates when it confidently says
yes to an object that isn't there. The intuition is that if visual neurons genuinely
carry visual evidence, then deleting them should leave the language model ``blind
but talkative'' --- it loses grounding in the image and hallucinates more. As a
mirror-image control, the same experiment is run on TriviaQA, a purely textual
general-knowledge quiz with no images at all. There, visual neurons should be
irrelevant, and instead deleting \emph{text} neurons should be what hurts. This
two-sided design is the heart of the experiment: visual damage should implicate
visual neurons, textual damage should implicate text neurons, and getting both
directions right is much harder to achieve by accident than either alone.

\section*{Step one: building a trustworthy question set}

Before measuring anything, the experiment cleans its measuring instrument. Models
answer stochastically, so a question the model gets right one time and wrong the
next would add noise to every measurement. Each candidate question is therefore
asked ten times with sampling turned on, and only the questions where the model is
\emph{consistent} --- always right or always wrong --- are kept. This is done across
all three POPE difficulty variants and capped to roughly a thousand questions, and
the same filtering is applied to TriviaQA to get about a thousand clean trivia
questions. The result is a stable behavioral baseline: we know exactly how the
intact model performs on every retained question.

\section*{Step two: silencing neurons and measuring the damage}

Now the slow part. The model's neurons are silenced in small groups --- fifty at a
time --- by clamping their output to zero during the forward pass, while everything
else in the model stays untouched. With one group silenced, the model re-answers the
entire thousand-question POPE set, and its hallucination rate is compared to the
intact baseline. The difference --- call it delta-H --- is the causal footprint of
that group: a positive delta-H means those neurons were \emph{suppressing}
hallucination, so removing them made things worse. The same silenced model also
re-answers the thousand TriviaQA questions, producing a second score for textual
damage. Then the hooks are removed, the next fifty neurons are silenced, and the
whole evaluation repeats. Since each group's score requires two thousand fresh model
generations, and a single layer contains hundreds of such groups, scoring one layer
takes most of a day on one GPU --- which is why the cluster runs one job per layer
and why this phase is where all the runs have been stuck. Each group's score is
attributed equally to the fifty neurons inside it; that's an acknowledged
approximation traded for tractability, since silencing neurons truly one at a time
would be fifty times slower.

\section*{Step three: the enrichment test itself}

Once every layer is scored, all neurons in the model are ranked by their delta-H,
and the top five percent are declared the ``hallucination-driving'' set --- the
neurons whose removal most increased hallucination. Now comes the statistical
question the experiment is named after: is any label category
\emph{over-represented} in that driving set relative to its share of the whole
model? Suppose visual neurons make up eight percent of all neurons but eighteen
percent of the driving set --- that's enrichment. Formally, for each category a
two-by-two table is built --- driving versus not driving, in the category versus not
--- and Fisher's exact test gives two things: an odds ratio (how many times more
likely a driving neuron is to carry that label than chance would predict, where
above one means enriched and below one means depleted) and a p-value for whether the
association could be a fluke. To calibrate everything, the same computation is also
run a thousand times on \emph{randomly chosen} neuron sets of the same size; these
random trials should show no enrichment, providing both a sanity check and bootstrap
confidence intervals. The analysis is also repeated layer by layer, producing a map
of where in the network's depth the enrichment concentrates.

\section*{What success looks like}

The hypothesis is a double dissociation. On the POPE side, the visual category
should come out significantly enriched among hallucination-driving neurons ---
deleting visual neurons blinds the model. On the TriviaQA side, the text category
should be the enriched one --- deleting text neurons erases knowledge. Multimodal
neurons may show up in both; unknown neurons and the random control should show
nothing. And because the whole experiment is run twice over --- once with the
permutation-test labels and once with the older fixed-threshold labels --- the paper
can make the comparative claim that matters: the statistically principled labels
predict causal function better than the heuristic ones. Every enrichment number
planned for the camera-ready comes out of exactly this machinery.

\end{document}
